\section*{Abstract}

High-performance computing programs spend most of their time in loop nests, and
the speed of the generated code depends on optimisation choices such as
vectorisation, unrolling and the use of the target instruction set. Production
compilers make these choices with fixed, hand-written heuristics that cannot
adapt to the individual loop. This report designs and evaluates an
ML-based loop-optimisation decision model that sits in front of an unmodified
GCC compiler.

A compiler front end is implemented from scratch: a lexical analyser, a
recursive-descent parser, a semantic and loop analyser, and a three-address-code
generator. From the analysed loop nest the front end computes \numFeatures\
static features, such as trip counts, exact memory strides, loop-carried
dependence distance, reductions, arithmetic intensity and working-set size. A
supervised model maps the features to one of \numConfigs\ GCC optimisation
configurations. The model is trained on \numKernels\ C loop kernels from
\numFamilies\ HPC-style families (vector updates, reductions, recurrences,
1-D/2-D stencils, matrix multiplication, 2-D traversals, branchy and strided
loops). Every kernel was compiled with every configuration (\numCompiles\
compilations) and executed on the target machine, and the label is the best
configuration that preserves the numerical result.

The evaluation uses the speed-up of the generated code. With random 5-fold
cross-validation the best model (\selModel) reaches a geometric-mean speed-up of
\selGMCV$\times$ over \texttt{-O2}, which is \selFracCV\% of the gain of the oracle
(\oracleGM$\times$). This must be compared honestly with static choices: always using
\texttt{-O3} gives \gmOThree$\times$, always using \texttt{-O3 -march=native} gives
\gmONative$\times$, and always using \texttt{-O3 -march=native -ffast-math} gives
\gmOFast$\times$. On this machine the learned model is therefore only marginally better
than the most aggressive static configuration in cross-validation, and on loop families
never seen in training (leave-one-family-out) it is worse (\selGMLOFO$\times$), which shows the
limits of a small corpus. The model is more useful when the binary must be portable and
IEEE-safe: among the four non-native, non-fast-math configurations it reaches \pCV$\times$
(oracle \pOracle$\times$) against \pOThree$\times$ for \texttt{-O3}, which causes \pOThreeSlow\
slow-downs of more than 5\%. Feature extraction and prediction take about \featMs\ ms and
\predMs\ ms per loop, small compared with a GCC compilation (\gccCompileMs\ ms).

\textbf{Keywords:}
Compiler Design, Loop Optimization, Lexical Analysis, Syntax Analysis,
Dependence Analysis, Machine Learning, High-Performance Computing, GCC.
